\documentclass[11pt, a4paper]{article}
\usepackage[utf8]{inputenc}
\usepackage{hyperref}
\usepackage{graphicx}
\usepackage{listings}
\usepackage{xcolor}
\usepackage{booktabs}
\usepackage{geometry}

\geometry{margin=1in}

\definecolor{codegray}{rgb}{0.5,0.5,0.5}
\definecolor{backcolour}{rgb}{0.95,0.95,0.92}

\lstdefinestyle{mystyle}{
    backgroundcolor=\color{backcolour},   
    commentstyle=\color{codegray},
    keywordstyle=\color{blue},
    numberstyle=\tiny\color{codegray},
    stringstyle=\color{teal},
    basicstyle=\ttfamily\footnotesize,
    breakatwhitespace=false,         
    breaklines=true,                 
    captionpos=b,                    
    keepspaces=true,                 
    numbers=left,                    
    numbersep=5pt,                  
    showspaces=false,                
    showstringspaces=false,
    showtabs=false,                  
    tabsize=2
}
\lstset{style=mystyle}

\title{\textbf{Advanced Hybrid-RAG Pipeline Architecture \& Developer Guide}}
\author{Pulkit Chauhan}
\date{\today}

\begin{document}

\maketitle

\tableofcontents
\newpage

\section{Introduction}
This document outlines the architecture, setup, and operational mechanics of the evolved Retrieval-Augmented Generation (RAG) system. Originally based on a standard semantic search prototype, the system has been fundamentally restructured to support hierarchical section-aware chunking, hybrid retrieval (BM25 + Cosine Vector Search), cross-encoder reranking, and decoupled configuration management.

\section{System Architecture}

The pipeline consists of the following core layers:
\begin{enumerate}
    \item \textbf{Ingestion Layer:} Handles batch PDF uploads, extracts structural elements, and pre-processes text.
    \item \textbf{Processing Layer (Chunking):} Implements Hierarchical Chunking to preserve contextual boundaries within documents.
    \item \textbf{Storage Layer:} Uses PostgreSQL extended with \texttt{pgvector} for storing embeddings and enriched metadata (page numbers, section headers).
    \item \textbf{Retrieval Layer:} A unified Hybrid Search engine combining exact-keyword matching (BM25) with semantic meaning (Cosine similarity).
    \item \textbf{Reranking Layer:} A precision bottleneck using local Cross-Encoder models (Jina Reranker v2) to elevate the most relevant chunks.
    \item \textbf{Generation Layer:} Supports both local LLMs (\texttt{llama-cpp-python}) ensuring data privacy, and ultra-fast cloud inference via the Groq API.
\end{enumerate}

\subsection{Implementation Variants (Architecture Diagrams)}

The repository contains four distinct execution scripts, each tailored for specific deployment profiles or testing scenarios.

\subsubsection{1. Unified Master Pipeline (\texttt{master\_pipeline.py})}
The new default flagship script. Instead of relying on separate files, this single entrypoint dynamically reconstructs its architecture at runtime based on the \texttt{pipeline\_strategy} defined in \texttt{master\_config.json}. It can seamlessly switch between Hierarchical and Standard chunking, Hybrid and Vector-only retrieval, and FlashRank, Jina, or No Reranking logic.

\begin{verbatim}
+-------------+      +-------------------+      +-------------------------+
| User (UI)   | ---> | FastAPI Backend   | ---> | 1. Config Routing       |
+-------------+      +-------------------+      +-------------------------+
                                                            |
                                                            v
+-------------------------+     +-------------------+     +-------------------------+
| 3. Dynamic Reranker     | <-- | 2. Dynamic Search | <-- | (Dynamic Chunking)      |
| (Jina/FlashRank/None)   |     | (Hybrid/Vector)   |     | (Hierarchical/Standard) |
+-------------------------+     +-------------------+     +-------------------------+
          |
          v
+-------------------------+     +-------------------+     +-------------+
| 4. System Prompt Builder| --->| 5. LLM Generation | --->| User (UI)   |
|    (Context + History)  |     |   (LlamaCpp Local)|     |  (Answer)   |
+-------------------------+     +-------------------+     +-------------+
\end{verbatim}

\subsubsection{2. Baseline Local RAG (\texttt{main\_local.py})}
+-------------+      +-------------------+      +-------------------------+
| User (UI)   | ---> | FastAPI Backend   | ---> | 1. Embed Query (MiniLM) |
+-------------+      +-------------------+      +-------------------------+
                                                            |
                                                            v
+-------------------------+     +-------------------+     +-------------------------+
| 3. Rerank (Jina Cross)  | <-- | 2. Vector Search  | <-- | (PostgreSQL pgvector)   |
+-------------------------+     +-------------------+     +-------------------------+
          |
          v
+-------------------------+     +-------------------+     +-------------+
| 4. System Prompt Builder| --->| 5. LLM Generation | --->| User (UI)   |
|    (Context + History)  |     |   (LlamaCpp Local)|     |  (Answer)   |
+-------------------------+     +-------------------+     +-------------+
\end{verbatim}

\subsubsection{2. Cloud Inference RAG (\texttt{main\_groq.py})}
Designed for ultra-low latency generation. It keeps embeddings and retrieval local for data privacy but offloads the final heavy LLM generation to the Groq Cloud API.

\begin{verbatim}
+-------------+      +-------------------+      +-------------------------+
| User (UI)   | ---> | FastAPI Backend   | ---> | 1. Embed Query (MiniLM) |
+-------------+      +-------------------+      +-------------------------+
                                                            |
                                                            v
+-------------------------+     +-------------------+     +-------------------------+
| 3. Rerank (Jina Cross)  | <-- | 2. Vector Search  | <-- | (PostgreSQL pgvector)   |
+-------------------------+     +-------------------+     +-------------------------+
          |
          v
+-------------------------+     +-------------------+     +-------------+
| 4. System Prompt Builder| --->| 5. LLM Generation | --->| User (UI)   |
|    (Context + History)  |     |  (Groq Cloud API) |     |  (Answer)   |
+-------------------------+     +-------------------+     +-------------+
\end{verbatim}

\subsubsection{3. Hierarchical + FlashRank (\texttt{main\_hierarchical\_flashrank.py})}
An intermediate script introducing Section-Aware (Hierarchical) Chunking. It swaps the heavy Jina model for the lighter FlashRank model for faster CPU-bound reranking.

\begin{verbatim}
+-------------+      +-------------------+      +-------------------------+
| User (UI)   | ---> | FastAPI Backend   | ---> | 1. Embed Query (MiniLM) |
+-------------+      +-------------------+      +-------------------------+
                                                            |
                                                            v
+-------------------------+     +-------------------+     +-------------------------+
| 3. Rerank (FlashRank)   | <-- | 2. Vector Search  | <-- | (pgvector + Sections)   |
+-------------------------+     +-------------------+     +-------------------------+
          |
          v
+-------------------------+     +-------------------+     +-------------+
| 4. System Prompt Builder| --->| 5. LLM Generation | --->| User (UI)   |
|    (Context + History)  |     |   (LlamaCpp Local)|     |  (Answer)   |
+-------------------------+     +-------------------+     +-------------+
\end{verbatim}

\subsubsection{4. Advanced Hybrid RAG (\texttt{main\_hierarchical\_hybrid.py})}
The most advanced pipeline. Combines Hierarchical chunking with a dual-retrieval system (Semantic Vector Search merged with BM25 Keyword Search), followed by heavy Jina reranking.

\begin{verbatim}
                                                          +-------------------------+
                                                  +-----> | Vector Search (pgvector)|
+-------------+      +-------------------+      +-+-+     +-------------------------+
| User (UI)   | ---> | FastAPI Backend   | ---> |1| Embed           | (Union)       |
+-------------+      +-------------------+      +-+-+     +-------------------------+
                                                  +-----> | BM25 Keyword Search     |
                                                          +-------------------------+
                                                                    |
                                                                    v
+-------------+      +-------------------+      +-------------------------+
| User (UI)   | <--- | 5. LLM Generation | <--- | 3. Rerank (Jina Cross)  |
|  (Answer)   |      |  (LlamaCpp Local) |      |                         |
+-------------+      +-------------------+      +-------------------------+
\end{verbatim}

\section{Core Components \& Configuration}

\subsection{Configuration Management (\texttt{master\_config.json})}
The system's behavior is entirely decoupled from the codebase via \texttt{master\_config.json}. This file acts as the ultimate control board, allowing operational adjustments, architectural routing, and hyperparameter tuning without requiring code modifications.

\begin{lstlisting}[language=json, caption=The Full master\_config.json File]
{
  "pipeline_strategy": {
    "chunking": "hierarchical",
    "search": "hybrid",
    "reranker": "jina"
  },
  "flush_database_on_shutdown": false,
  "llm_path": "models/general_model.gguf",
  "embedding_model": "all-MiniLM-L6-v2",
  "chunking_tokens": {
    "max_size_default": 650,
    "max_size_list_item": 600,
    "max_size_short_def": 400,
    "overlap": 100
  },
  "llm_hyperparameters": {
    "temperature": 0.0,
    "max_tokens": 2048,
    "n_ctx": 4096,
    "n_gpu_layers": -1,
    "n_batch": 512,
    "n_threads": 8,
    "repeat_penalty": 1.1,
    "stop": ["<|eot_id|>", "<|start_header_id|>", "Human:", "Assistant:"],
    "verbose": false
  },
  "retriever_hyperparameters": {
    "top_k": 20,
    "top_n_rerank": 3,
    "rrf_k": 60,
    "base_similarity_limit": 5
  },
  "llm_chat_template": {
    "contextualize_q_system_prompt": "Given a chat history and the latest user...",
    "qa_system_prompt": "You are a helpful assistant for the platform..."
  }
}
\end{lstlisting}

\subsubsection{Key configuration parameters explained:}
\begin{itemize}
    \item \textbf{\texttt{pipeline\_strategy}:} The dynamic router. 
    \begin{itemize}
        \item \texttt{chunking}: Set to \texttt{"hierarchical"} to split PDFs by Section Headings, or \texttt{"standard"} for recursive cuts.
        \item \texttt{search}: Set to \texttt{"hybrid"} to pool Vector L2 distance with BM25 Keyword Matching, or \texttt{"vector"} to skip BM25 entirely (saves RAM).
        \item \texttt{reranker}: Set to \texttt{"jina"} for heavy precise cross-encoding, \texttt{"flashrank"} for fast CPU rescoring, or \texttt{"none"} to pass chunks directly.
    \end{itemize}
    \item \textbf{\texttt{flush\_database\_on\_shutdown}:} If true, wiping the PostgreSQL records upon server termination.
    \item \textbf{\texttt{chunking\_tokens.max\_size\_default}:} Dictates the standard size of text chunks. Larger sizes provide more context to the LLM but consume more tokens; smaller sizes increase retrieval precision.
    \item \textbf{\texttt{llm\_hyperparameters}:} Deep tuning for local \texttt{llama.cpp} execution. \texttt{n\_ctx} determines your context window size. \texttt{n\_gpu\_layers = -1} offloads all possible layers to hardware acceleration.
    \item \textbf{\texttt{retriever\_hyperparameters.top\_k}:} The initial pool of chunks retrieved from the database before reranking.
    \item \textbf{\texttt{retriever\_hyperparameters.top\_n\_rerank}:} The final number of high-precision chunks that survive the reranker and are passed strictly to the LLM as context.
\end{itemize}

\subsection{Data Ingestion \& Hierarchical Chunking}
Standard RAG systems blindly chop text into n-character blocks, often splitting paragraphs in half. 

\textbf{Our Implementation:} The \texttt{main\_hierarchical\_hybrid.py} script employs Regex-based structure parsing. It identifies Chapter and Section headers within the PDF text. When chunking text beneath a header, the chunk is prefixed with metadata (e.g., \textit{"Section: 3.1 Architecture Overview"}). 

\textit{Why this matters:} When the LLM receives an isolated snippet of text retrieved from the vector database, the embedded section header immediately informs the model of the text's original context within the larger document.

\subsection{Hybrid Retrieval Mechanism}
Standard Vector Databases struggle with exact keyword lookups (e.g., finding an exact ID number "SYS-9941").

\textbf{Our Implementation:}
\begin{enumerate}
    \item \textbf{Vector Search:} Calculates Cosine Distance between the user's query embedding and the chunk embeddings via PostgreSQL's \texttt{pgvector}. This captures the "meaning" of the question.
    \item \textbf{BM25 (Best Matching 25):} A local, in-memory inverted index that tracks term frequency. It is excellent at exact keyword matching.
    \item \textbf{Union \& Deduplication:} The Top $K$ results from Vector Search and the Top $K$ results from BM25 are pooled together into a single unique set of candidate chunks.
\end{enumerate}

\subsection{Reranking via Jina Cross-Encoder}
The combined candidate pool from the Hybrid Retreiver (often 20-30 chunks) is too noisy to send directly to an LLM. 

\textbf{Our Implementation:} The system evaluates every candidate chunk against the user's query using the \texttt{Jina Reranker v2} model (loaded locally via GGUF). Cross-encoders process the query and the chunk simultaneously, calculating a highly accurate relevance score. Only the highest-scoring chunks (\texttt{top\_n\_rerank}) survive this phase.

\section{Language Models: Local vs. Cloud}

\subsection{Local Inference (\texttt{main\_hierarchical\_hybrid.py})}
Uses \texttt{llama-cpp-python} to run models (like Llama 3) entirely locally.
\begin{itemize}
    \item \textbf{Pros:} Absolute data privacy; no API costs; no internet connection required.
    \item \textbf{Cons:} Slower inference speeds bound by local CPU/GPU hardware limitations.
\end{itemize}

\subsection{Cloud Inference (\texttt{main\_groq.py})}
Uses the Groq API (e.g., \texttt{llama-3-70b-8192}).
\begin{itemize}
    \item \textbf{Pros:} Extremely low latency (tokens per second); ability to use massive 70B+ parameter models.
    \item \textbf{Cons:} Requires an internet connection and an active API Key (\texttt{GROQ\_API\_KEY}). Data is sent to Groq's servers for inference.
\end{itemize}

\section{User Interfaces \& Frontends}

\subsection{Streamlit Web Interface (\texttt{frontend.py})}
The repository includes a responsive web application built with Streamlit.

Features include:
\begin{itemize}
    \item \textbf{Multi-File Uploading:} Users can drag and drop multiple PDF files simultaneously. The UI holds a loading state while the backend performs hierarchical chunking and embedding.
    \item \textbf{Session Isolation:} Utilizes \texttt{uuid4()} to generate continuous conversational history specific to the user's browser tab.
    \item \textbf{Source Attribution:} The chat UI explicitly displays the retrieved sources (document names and confidence scores) underneath every generated answer for transparency.
    \item \textbf{Database State Management:} Provides a secure ``Clear Chat History'' button that also communicates with the backend to flush the database via the \texttt{/clear-db/} endpoint.
\end{itemize}

\section{Database Maintenance \& Operations}

\subsection{Schema \& Setup}
Run \texttt{python create\_db.py} to initialize the PostgreSQL database. This script ensures the required \texttt{pgvector} extension is active and builds the necessary \texttt{TextChunk} tables defined in \texttt{models.py}.

\subsection{Flushing \& Cleanup}
During testing, the database can become bloated.
\begin{itemize}
    \item \textbf{Manual Utility:} Run \texttt{python flush\_db.py} to safely wipe all \texttt{TextChunk} records and delete the physical \texttt{bm25*.pkl} index files from the disk.
    \item \textbf{Automated Shutdown Hook:} If \texttt{flush\_database\_on\_shutdown} is true in \texttt{config.json}, the FastAPI server will automatically clear the database on termination.
    \item \textbf{Frontend UI:} The Streamlit app includes a "Clear Chat History" button that invokes a dedicated API endpoint to wipe the backend state.
\end{itemize}

\section{Automated Testing \& Observation}

\subsection{Robust Logging}
Every successful interaction via the \texttt{/chat/} endpoint generates a comprehensive JSON log within the \texttt{chat history/} directory. 
This log contains:
\begin{itemize}
    \item The original user query.
    \item The exact text chunks retrieved.
    \item The reranking score assigned to each chunk (allowing developers to tune similarity limits).
    \item The fully constructed LLM prompt.
    \item Latency (in milliseconds) for the full request lifecycle.
\end{itemize}

\subsection{RAG Evaluation \& Stability Testing}

To systematically evaluate the performance of the RAG pipeline over time, the repository includes an automated testing framework: \texttt{test\_model\_with\_restart.py}.

This framework is built to mirror real-world usage safely, accounting for long-running LLM processes that might leak memory.

\subsubsection{How the Testing Framework Operates}
When executed, the script orchestrates the following lifecycle:
\begin{enumerate}
    \item \textbf{Backend Initialization:} Spawns the \texttt{main\_hierarchical\_hybrid.py} server as a subprocess and continuously polls the \texttt{/docs} endpoint until the server is strictly ready.
    \item \textbf{Automated Ingestion:} Posts the designated test document (\texttt{evaluation/input.pdf}) to the \texttt{/upload-pdfs/} endpoint to populate the fresh database.
    \item \textbf{Batch Execution:} Reads a series of evaluation queries line-by-line from a CSV file (\texttt{rag queries - Sheet1.csv}) and sends them to the \texttt{/chat/} endpoint.
    \item \textbf{Preemptive Restarts:} After a configurable number of queries (\texttt{QUERIES\_PER\_RESTART}, default is 3), the script violently terminates the backend subprocess, waits for ports to clear, and respawns the server. This guarantees that the system is tested under clean-state conditions and proves resilience against memory bloating over hundreds of requests.
\end{enumerate}

\subsubsection{Usage Instructions}

To run an evaluation batch:
\begin{enumerate}
    \item Ensure your virtual environment is active.
    \item Place the target document you wish to test into the \texttt{evaluation/} folder and name it \texttt{input.pdf}.
    \item Populate your questions in \texttt{rag queries - Sheet1.csv} (one question per row).
    \item Execute the script from the terminal:
\end{enumerate}

\begin{lstlisting}[language=bash, caption=Running the Evaluation Suite]
python test_model_with_restart.py
\end{lstlisting}

Monitor the console output. The script will detail the server start phases, document upload success, and print a snippet of every generated answer as the batch progresses. All detailed logs will still be written cleanly out to the \texttt{chat history/} folder.

\section{Dependencies \& Environment Setup}

\textbf{CRITICAL REQUIREMENT:} You MUST be running \textbf{Python 3.14.2} strictly for this environment to function correctly. Differences in Python versions may result in unresolvable native binary conflicts with \texttt{llama-cpp-python} or \texttt{pgvector}.

The environment dependencies ensure compatibility across the ML stack, vector database integration, and UI rendering. The following is the updated \texttt{requirements.txt} mapping generated via \texttt{pip freeze}:

\begin{lstlisting}[language=bash, caption=requirements.txt details]
boto3==1.42.63
botocore==1.42.63
fastapi==0.135.1
FlashRank==0.2.10
huggingface_hub==1.4.1
langchain_classic==1.0.2
langchain_community==0.4.1
langchain_core==1.2.17
langchain_groq==1.1.2
langchain_openai==1.1.10
langchain_text_splitters==1.1.1
llama_cpp_python
pgvector==0.4.2
psycopg2_binary==2.9.11
pydantic==2.12.5
pypdf==6.8.0
python-dotenv==1.2.2
requests==2.32.5
sentence_transformers==5.2.3
SQLAlchemy==2.0.48
streamlit==1.54.0
uvicorn==0.41.0
\end{lstlisting}

\textbf{To install:}
\begin{lstlisting}[language=bash]
pip install -r requirements.txt
\end{lstlisting}

\end{document}
